% !Mode:: "TeX:UTF-8"
% !Mode:: "TeX:UTF-8"
\newcommand{\thepage}{} % to silence a hyperref warning about pdfpagelabels
\RequirePackage[]{silence}
\WarningsOff[hyperref]
\PassOptionsToPackage{log-declarations=false}{xparse}
%\RequirePackage[l2tabu, orthodox]{nag}
\RequirePackage{luatex85} % to make \pdf...= commands work with LuaTeX
\documentclass[aspectratio=169,t]{beamer}
\usepackage{ifthen}
\usepackage[T1]{fontenc}
\usepackage{iftex}
\ifLuaTeX
 \usepackage{fontspec}
 \setmainfont[Ligatures=TeX]{Source Sans Pro}
\else
 \usepackage[utf8]{inputenc}
\fi
\usepackage{german}
\usepackage{newunicodechar}
\newunicodechar{€}{\texteuro}
\newunicodechar{ß}{\ss}
\renewcommand{\SS}{\iffontchar\font"1E9E \symbol{"1E9E}\else SS\fi}
\usepackage{ellipsis}
\usepackage{microtype}
\usepackage{listings}
\usepackage{sourcecodepro}
\lstset{language=Python, basicstyle=\ttfamily}
\usepackage{amsfonts}
\usepackage{amssymb}
\usepackage{amsthm}
\usepackage{fancyhdr}
\usepackage{graphicx}

\newtheorem{proposition}[theorem]{Proposition}

\DeclareMathOperator*{\argmax}{argmax}
\DeclareMathOperator{\supp}{supp}
\newcommand{\E}{\ensuremath{\mathbb{E}}}
\newcommand{\V}{\ensuremath{\mathbb{V}}}

\setbeamercolor{emph}{fg=purple}
\renewcommand<>{\emph}[1]{%
  {\usebeamercolor[fg]{emph}\only#2{\itshape}#1}%
}



\hypersetup{draft=false,colorlinks=false,pdfpagemode=UseOutlines,plainpages=false,
            breaklinks=true,bookmarksnumbered=true,pdfborder=0 0 0,
            pdftitle={Einführung in Deep Learning WiSe 2026/27 HHU},
            pdfauthor={Konrad Voelkel},
            pdfsubject={data science}}

%\documentclass[aspectratio=169,t]{beamer}
\usetheme{hhubeamer} % put the hhu-beamer-theme in /usr/share/texmf/tex/latex/beamer/themes/

% -------------------- Code listings config --------------------
% Define custom colors for syntax highlighting
\definecolor{codegreen}{rgb}{0,0.6,0}
\definecolor{codegray}{rgb}{0.5,0.5,0.5}
\definecolor{codepurple}{rgb}{0.58,0,0.82}
\definecolor{backcolour}{rgb}{0.96,0.96,0.94} % Very light gray/beige
% Configure the listings package
\lstset{
    language=Python,
    backgroundcolor=\color{backcolour},   
    commentstyle=\color{codegreen}\itshape,
    keywordstyle=\color{blue}\bfseries,
    numberstyle=\tiny\color{codegray},
    stringstyle=\color{codepurple},
    basicstyle=\ttfamily\footnotesize, % Use a monospaced font
    breakatwhitespace=false,         
    breaklines=true,      
             % Enable automatic line breaking
    captionpos=b,                      % Caption at the bottom
    keepspaces=true,                 
    numbers=left,                      % Line numbers on the left
    numbersep=10pt,                    % Distance of line numbers from the code
    showspaces=false,                
    showstringspaces=false,            % Don't show weird underscores in strings
    showtabs=false,                  
    tabsize=4,
    frame=single,                      % Add a subtle border around the code
    rulecolor=\color{black!20},        % Border color
  }

\event[]{}
\author[Wintersemester 2026/27]{Dr. Konrad Völkel, Wintersemester 2026/27}

%\title{Data Science, VL 1}
%\date{16.04.2026}


%%% Local Variables:
%%% mode: latex
%%% TeX-master: t
%%% End:


\title{Modulvorstellung \\Einführung in Deep Learning (EiDL)}
\date{31.07.2026}


\begin{document}

\begin{frame}
  \titlepage
\end{frame}

\begin{frame}
  \frametitle{Voraussetzungen, Zulassung, Prüfung}

  \emph{Formale Voraussetzungen}: keine

  \emph{Inhaltliche Voraussetzungen}:  Inhalte der Module \textbf{Data Science}, denn wir werden rudimentäre Python-Kenntnisse und Grundlagen des maschinellen Lernens sowie insb. die vektorisierte Programmierung mit NumPy und die Idee der linearen Regression voraussetzen.

  \emph{Zulassung zur Prüfung}: Alle (mindestens 10, höchstens 12) Übungsblätter müssen bearbeitet werden und mindestens 8 abgegeben und als ``bestanden'' ( = mind. die Hälfte der Punkte ) bewertet worden sein.

  \emph{Prüfungsform}: schriftlich (Klausur, i.d.R. 90 Minuten) ohne Hilfsmittel/Computer,\\
  Fokus auf Theorie.

  Die \emph{Note} wird vollständig aus der Klausur gebildet, die Übungsblätter zuvor sind lediglich eine Studienleistung.
  
\end{frame}

\begin{frame}
  \frametitle{Material zur Übung- und Prüfungsvorbereitung}

  \begin{itemize}
  \item Vorlesung (Theorie) immer mittwochs 12:30 in Hörsaal 5E.
    \item Übungen (Besprechung der Musterlösung, Fragestunde und zusätzliche Anwesenheitsaufgaben unter Anleitung bearbeiten) Do 10:30 in 2522.U1.52.
\item  Aus 2025/26 gibt es VL-Aufzeichnungen in der Mediathek,\\
  \url{media.hhu.de/search/title/Einführung+in+Deep+Learning+2025/}
\item  Die Übungsaufgaben kommen als wöchentliches PDF,\\größtenteils mit Musterlösung.
\item  Die Vorlesungsfolien sind als PDF wöchentlich vorab verfügbar.
\item Ein Skript gibt es darüber hinaus nicht.
\end{itemize}
\end{frame}

\begin{frame}
  \frametitle{Inhalt}
  \framesubtitle{Grober zeitlicher Ablauf}

  \begin{enumerate}
  \item Geschichte (Perzeptron)
    \item Logistische Regression ($\sim$ 1 Neuron) ausführlichst
    \item Stochastischer Gradientenabstieg
    \item Backpropagation: effiziente Berechnung der Ableitung für tiefe Netze
    \item \texttt{PyTorch} für Inferenz und Training from Scratch
    \item Autoencoder und VAE (unsupervised Representation Learning)
    \item Latente Vektoren für Natural Language Processing (NLP) mit word2vec
    \item Intro Convolutional Neural Networks (CNN)
    \item Intro Generative-Adversarial Networks (GAN)
    \item Vortrainierte LLMs in DL-Anwendungen nutzen
    \item Einzelne Exkurse/Vertiefungen nach Laune
  \end{enumerate}
  
\end{frame}


\begin{frame}
  \frametitle{Was wir nicht machen werden}

  \begin{itemize}
  \item Transformer-Modelle trainieren (GPT usw.)
  \item Auf großen echten Daten arbeiten
  \item Aktuelle Architekturen nachimplementieren
  \item Optimierungstheorie
  \item CUDA-Programmierung
  \item Reinforcement Learning
  \end{itemize}
\end{frame}


\begin{frame}
  \frametitle{Inhalt (Modulhandbuch)}
  \framesubtitle{Nach erfolgreicher Teilnahme können die Studierenden:}

\begin{itemize}
\item Grundbegriffe und -konzepte neuronaler Netze benennen und erklären,
\item  mathematische Grundlagen für neuronale Netze anwenden,
\item  einfache Modelle selbst implementieren,
\item  vortrainierte Modelle in andere Systeme integrieren und
\item  beurteilen, welche der besprochenen Modelle hinsichtlich einer Anwendung in Frage kommen.
\end{itemize}

\end{frame}

\begin{frame}
  \frametitle{Relation zu anderen Modulen in Bachelor/Master}

  
  \begin{itemize}
    \item Wenig bis keine Überschneidung:\\
      Machine Learning, Data Science 2, Introduction to NLP, Python for NLP
      \item Vertiefung:\\
        Deep Learning im Master, zahlreiche darauf aufbauende Kurse wie etwa\\
        Advanced Deep Learning, Implementing Transformer Models, Representations in Natural Language Processing, Probabilistic Machine Learning, Reinforcement Learning
  \end{itemize}
  
\end{frame}



\end{document}

%%% Local Variables:
%%% mode: latex
%%% TeX-master: t
%%% End:
